\appendix

% ============================================================================
\section{Experimental Setup Details} \label{app:setup}
% ============================================================================

\paragraph{Datasets and splits.}
We evaluate on four domains presented sequentially: 
GSM8K \citep{cobbe2021gsm8k} (\texttt{train} split), 
MMLU \citep{hendrycks2021mmlu} (\texttt{test} split), 
ARC-Challenge \citep{clark2018arc} (\texttt{test} split), 
and TruthfulQA \citep{lin2022truthfulqa} (\texttt{validation} split, MC1 variant).
We sample 500 questions per domain for a total of 2{,}000 questions per run.
The default domain order is GSM8K $\to$ MMLU $\to$ ARC $\to$ TruthfulQA.

\paragraph{Models.}
We evaluate on three instruction-tuned LLMs spanning 2--4B parameters:
Llama~3.2-3B-Instruct \citep{touvron2023llama},
Gemma~2-2B-IT \citep{team2024gemma},
and Phi~3.5-Mini-Instruct \citep{abdin2024phi}.

\paragraph{Trainable parameters.}
Table~\ref{tab:lora_params} reports the exact LoRA parameter counts per model.
The variation reflects architectural differences (Phi uses a fused \texttt{qkv\_proj} 
module rather than separate \texttt{q\_proj}/\texttt{v\_proj}).

\begin{table}[h]
\centering
\caption{LoRA trainable parameters per model (rank $r{=}8$).
Llama targets the last 4 layers (layers 24--27 of 28);
Gemma and Phi target the last 8 layers 
(Gemma: layers 18--25 of 26; Phi: layers 24--31 of 32).}
\label{tab:lora_params}
\small
\begin{tabular}{@{}lcccc@{}}
\toprule
\textbf{Model} & \textbf{Total params} & \textbf{LoRA layers} & \textbf{LoRA params} & \textbf{\% of total} \\
\midrule
Llama 3.2-3B  & 3.21B & 4  & 327{,}680 & 0.010\% \\
Gemma 2-2B    & 2.61B & 8  & 491{,}520 & 0.019\% \\
Phi 3.5-Mini  & 3.82B & 8  & 786{,}432 & 0.021\% \\
\bottomrule
\end{tabular}
\end{table}


% ============================================================================
\section{Hyperparameter Settings} \label{app:hyperparams}
% ============================================================================

\begin{table}[h]
\centering
\caption{Full hyperparameter settings used across all experiments.
Model-specific values are noted where applicable.}
\label{tab:hyperparams}
\small
\begin{tabular}{@{}llc@{}}
\toprule
\textbf{Component} & \textbf{Parameter} & \textbf{Value} \\
\midrule
\multirow{6}{*}{LoRA}
  & Rank $r$ & 8 \\
  & Scaling factor $\alpha$ & 16 \\
  & Target layers (Llama) & Last 4 \\
  & Target layers (Gemma, Phi) & Last 8 \\
  & Target modules (Llama, Gemma) & \texttt{q\_proj}, \texttt{v\_proj} \\
  & Target modules (Phi) & \texttt{qkv\_proj} \\
\addlinespace
\multirow{3}{*}{Optimization}
  & Optimizer & AdamW \\
  & Learning rate & $5 \times 10^{-5}$ \\
  & Epochs per question & 3 \\
\addlinespace
\multirow{2}{*}{Directional loss}
  & Step size $\alpha_{\text{step}}$ & 0.5 \\
  & Clip bound $\delta$ & 0.15 \\
\addlinespace
\multirow{1}{*}{Bin-gate}
  & Threshold & 1 bin \\
\addlinespace
\multirow{4}{*}{PH detector}
  & Tolerance $\epsilon$ & 0.05 \\
  & Detection threshold $\lambda$ & 3.0 \\
  & EMA smoothing $\alpha_{\text{ema}}$ & 0.05 \\
  & Warmup period & 30 questions \\
\addlinespace
\multirow{1}{*}{Burst}
  & Burst size $B$ & 50 \\
\addlinespace
\multirow{3}{*}{Normalization $\tau$}
  & Llama 3.2-3B & 0.7 \\
  & Gemma 2-2B & 1.5 \\
  & Phi 3.5-Mini & 1.5 \\
\addlinespace
\multirow{2}{*}{Other}
  & Weight accumulation & On \\
  & Distractors $k$ & 4 \\
\bottomrule
\end{tabular}
\end{table}


% ============================================================================
\section{Main Results} \label{app:main_results}
% ============================================================================

Table~\ref{tab:main_results} reports overall metrics for all methods.
We include both standard ECE (10 equal-width bins) and Adaptive ECE
(AdaECE; 10 equal-mass bins) to verify that improvements are robust
to the choice of binning strategy.

\begin{table}[h]
\centering
\caption{Overall results across all models and methods (2{,}000 questions each).
\textbf{Bold}: best overall per model per metric.
\underline{Underline}: best among label-free methods.}
\label{tab:main_results}
\small
\begin{tabular}{@{}llccccc@{}}
\toprule
\textbf{Model} & \textbf{Method}
  & \textbf{ECE}$\downarrow$ & \textbf{AdaECE}$\downarrow$ & \textbf{Brier}$\downarrow$
  & \textbf{AUROC}$\uparrow$ & \textbf{Acc} \\
\midrule
\multirow{6}{*}{Llama 3.2-3B}
  & Soft Verbalized     & .170 & .167 & .292 & .510 & .576 \\
  & Self-Consistency    & .111 & --- & .215 & \textbf{.718} & .624 \\
  & ~~+ Temp Scaling$^*$ & .047 & .075 & .250 & .504 & .576 \\
  & P(True) Norm        & .065 & .064 & .223 & .693 & .576 \\
  & ~~+ Temp Scaling$^*$ & \textbf{.032} & \textbf{.035} & \textbf{.222} & .676 & .576 \\
  & SECL (Ours)         & \underline{.050} & \underline{.060} & .241 & .587 & .577 \\
\addlinespace
\multirow{5}{*}{Gemma 2-2B}
  & Soft Verbalized     & .256 & .252 & .314 & .558 & .516 \\
  & ~~+ Temp Scaling$^*$ & .047 & .037 & .249 & .550 & .516 \\
  & P(True) Norm        & .141 & .144 & .258 & \textbf{.650} & .516 \\
  & ~~+ Temp Scaling$^*$ & \textbf{.037} & \textbf{.033} & \textbf{.233} & .647 & .516 \\
  & SECL (Ours)         & \underline{.056} & \underline{.060} & .254 & .548 & .515 \\
\addlinespace
\multirow{5}{*}{Phi 3.5-Mini}
  & Soft Verbalized     & .251 & .240 & .275 & .600 & .667 \\
  & ~~+ Temp Scaling$^*$ & \textbf{.047} & \textbf{.049} & .215 & .583 & .667 \\
  & P(True) Norm        & .154 & .151 & .227 & \textbf{.675} & .667 \\
  & ~~+ Temp Scaling$^*$ & .059 & .060 & \textbf{.205} & .664 & .667 \\
  & SECL (Ours)         & \underline{.115} & \underline{.119} & .251 & .521 & .665 \\
\bottomrule
\end{tabular}
\vspace{0.3em}
\newline
{\footnotesize $^*$Requires ground-truth labels; fitted via 5-fold cross-validation. See \S\ref{app:posthoc}.}
\end{table}


% ============================================================================
\section{Post-Hoc Calibration Baselines} \label{app:posthoc}
% ============================================================================

We compare SECL against standard post-hoc calibration methods:
\textbf{temperature scaling} \citep{guo2017calibration} (single parameter $T$
fitted to minimize NLL) and \textbf{Platt scaling} \citep{platt1999probabilistic}
(logistic regression $\sigma(a \cdot c + b)$).
Both are applied to the soft verbalized confidence and fitted via 5-fold 
cross-validation (each fold uses 1{,}600 labeled examples for fitting).

\paragraph{Key difference from SECL.}
Temperature and Platt scaling are \emph{supervised} post-hoc methods:
they require ground-truth correctness labels to fit their parameters.
SECL is \emph{unsupervised} at test time---it uses only the model's own 
P(True) signal, with no access to ground truth.
Additionally, post-hoc methods are monotone transforms that 
\emph{cannot} improve discrimination (AUROC);
SECL can improve AUROC by adapting the model's representations.

\begin{table}[h]
\centering
\caption{Post-hoc calibration baselines (5-fold CV). 
Temperature scaling compresses the confidence range substantially,
achieving low ECE by predicting near-constant values for models with
weak verbalized discrimination (fitted $T$ shown).}
\label{tab:posthoc}
\small
\begin{tabular}{@{}llccccc@{}}
\toprule
\textbf{Model} & \textbf{Method}
  & \textbf{ECE}$\downarrow$ & \textbf{AdaECE}$\downarrow$
  & \textbf{Brier}$\downarrow$ & \textbf{AUROC}$\uparrow$ 
  & \textbf{Conf.\ range} \\
\midrule
\multirow{4}{*}{Llama}
  & SV + Temp ($T{=}17.6$)  & .047 & .075 & .250 & .504 & [.44,\,.56] \\
  & SV + Platt              & .021 & .057 & .244 & .481 & [.56,\,.61] \\
  & PTN + Temp ($T{=}1.8$)  & .032 & .035 & .222 & .676 & [.15,\,.85] \\
  & SECL (no labels)        & .050 & .060 & .241 & .587 & [.05,\,.85] \\
\addlinespace
\multirow{4}{*}{Gemma}
  & SV + Temp ($T{=}11.2$)  & .047 & .037 & .249 & .550 & [.41,\,.59] \\
  & SV + Platt              & .035 & .050 & .250 & .550 & [.42,\,.55] \\
  & PTN + Temp ($T{=}3.0$)  & .037 & .033 & .233 & .647 & [.26,\,.74] \\
  & SECL (no labels)        & .056 & .060 & .254 & .548 & [.05,\,.95] \\
\addlinespace
\multirow{4}{*}{Phi}
  & SV + Temp ($T{=}3.4$)   & .047 & .049 & .215 & .583 & [.29,\,.71] \\
  & SV + Platt              & .052 & .050 & .216 & .585 & [.16,\,.69] \\
  & PTN + Temp ($T{=}2.4$)  & .059 & .060 & .205 & .664 & [.21,\,.79] \\
  & SECL (no labels)        & .115 & .119 & .251 & .521 & [.05,\,.95] \\
\bottomrule
\end{tabular}
\end{table}

\paragraph{Discussion.}
For Llama, temperature scaling on the soft verbalized confidence requires 
$T{=}17.6$, compressing all predictions to [0.44,\,0.56]---essentially 
predicting the base rate.
Platt scaling is more extreme: its output range [0.56,\,0.61] is near-constant,
yielding ECE$=$0.021 but \emph{degrading} AUROC to 0.481 (below chance).
These methods achieve low ECE by destroying discriminative information.
By contrast, SECL achieves ECE$=$0.050 while \emph{improving} AUROC
from 0.510 to 0.587 (+0.077), demonstrating genuine calibration rather
than confidence compression.

For Phi, where the verbalized confidence carries more signal 
(baseline AUROC$=$0.600), temperature scaling uses a more moderate 
$T{=}3.4$ and achieves genuinely good calibration (ECE$=$0.047).
The strongest supervised baseline across all models is PTN~+~Temp,
which combines the discriminative P(True) signal with post-hoc recalibration.

\paragraph{Complementarity: SECL + Temperature Scaling.}
SECL and post-hoc calibration are complementary: applying temperature 
scaling to SECL's output further reduces ECE, since SECL improves the 
underlying confidence signal while temperature scaling recalibrates it.

\begin{table}[h]
\centering
\caption{SECL combined with temperature scaling (5-fold CV).
SECL+Temp achieves the best or near-best ECE overall while preserving
SECL's discrimination improvements.}
\label{tab:secl_plus_temp}
\small
\begin{tabular}{@{}llcccc@{}}
\toprule
\textbf{Model} & \textbf{Method}
  & \textbf{ECE}$\downarrow$ & \textbf{AdaECE}$\downarrow$
  & \textbf{Brier}$\downarrow$ & \textbf{AUROC}$\uparrow$ \\
\midrule
\multirow{3}{*}{Llama}
  & PTN + Temp$^*$  & .032 & .035 & .222 & .676 \\
  & SECL            & .050 & .060 & .241 & .587 \\
  & SECL + Temp$^*$ & .049 & .053 & .241 & .584 \\
\addlinespace
\multirow{3}{*}{Gemma}
  & PTN + Temp$^*$  & .037 & .033 & .233 & .647 \\
  & SECL            & .056 & .060 & .254 & .548 \\
  & SECL + Temp$^*$ & .011 & .037 & .249 & .542 \\
\addlinespace
\multirow{3}{*}{Phi}
  & PTN + Temp$^*$  & .059 & .060 & .205 & .664 \\
  & SECL            & .115 & .119 & .251 & .521 \\
  & SECL + Temp$^*$ & .097 & .082 & .232 & .516 \\
\bottomrule
\end{tabular}
\vspace{0.3em}
\newline
{\footnotesize $^*$Requires ground-truth labels for temperature fitting.}
\end{table}

\paragraph{Calibration--discrimination trade-off.}
SECL's primary objective is calibration (ECE), not discrimination (AUROC).
Table~\ref{tab:main_results} reveals a model-dependent trade-off:
for Llama, SECL \emph{improves} AUROC from .510 to .587 (+.077)
while reducing ECE from .170 to .050,
indicating that the adapted confidence carries genuinely better signal.
For Gemma, AUROC is roughly preserved (.558 $\to$ .548, $-$.010)
despite a large ECE reduction (.256 $\to$ .056).
For Phi, the trade-off is more pronounced:
AUROC drops from .600 to .521 ($-$.079) as ECE improves from .251 to .115.
Crucially, Brier score---which jointly captures calibration and
discrimination---improves for all three models
(Llama: .292 $\to$ .241; Gemma: .314 $\to$ .254; Phi: .275 $\to$ .251),
confirming that the net effect of SECL is beneficial.
By comparison, supervised temperature scaling also reduces AUROC
(Llama: $-$.006; Gemma: $-$.008; Phi: $-$.017)
and Platt scaling degrades it further for Llama ($-$.029, below chance at .481),
demonstrating that AUROC loss is not unique to SECL but is an inherent
cost of recalibrating overconfident models.


% ============================================================================
\section{Hyperparameter Sensitivity} \label{app:hp_sensitivity}
% ============================================================================

Table~\ref{tab:hp_sensitivity} reports ECE and Adaptive ECE on 
Llama~3.2-3B when varying individual hyperparameters.
The $\alpha_{\text{step}}$ and $\delta$ ablations use the full pipeline 
\emph{without} PH-gating (all questions trained) to isolate the loss 
function's effect; the burst size $B$ ablation uses the PH-gated pipeline.

\begin{table}[h]
\centering
\caption{Hyperparameter sensitivity (Llama~3.2-3B).
Bold indicates the selected configuration.
$^\dagger$Without PH-gating.
$^\ddagger$With PH-gating.}
\label{tab:hp_sensitivity}
\small
\begin{tabular}{@{}llcc@{}}
\toprule
\textbf{Parameter} & \textbf{Value} & \textbf{ECE}$\downarrow$ & \textbf{AdaECE}$\downarrow$ \\
\midrule
\multirow{3}{*}{$\alpha_{\text{step}}$$^\dagger$}
  & 0.2 & .067 & .083 \\
  & 0.3 & .066 & .069 \\
  & \textbf{0.5} & \textbf{.052} & \textbf{.073} \\
\addlinespace
\multirow{2}{*}{$\delta$$^\dagger$}
  & \textbf{0.15} & \textbf{.052} & \textbf{.073} \\
  & 0.20 & .064 & .067 \\
\addlinespace
\multirow{1}{*}{Loss$^\dagger$}
  & Plain MSE (no directional) & .085 & .088 \\
\addlinespace
\multirow{2}{*}{Burst $B$$^\ddagger$}
  & 20 & .114 & --- \\
  & \textbf{50} & \textbf{.050} & \textbf{.060} \\
\addlinespace
\multirow{3}{*}{$\tau$ (per model)}
  & Llama: \textbf{0.7} & \textbf{.065} & \textbf{.064} \\
  & Gemma: \textbf{1.5} & .141 & .144 \\
  & Phi: \textbf{1.5} & .154 & .151 \\
\bottomrule
\end{tabular}
\end{table}


% ============================================================================
\section{Per-Domain Results} \label{app:per_domain}
% ============================================================================

\begin{table}[h]
\centering
\caption{Per-domain results for Llama~3.2-3B.
Best ECE per domain in \textbf{bold}.}
\label{tab:per_domain_llama}
\small
\begin{tabular}{@{}llcccccc@{}}
\toprule
\textbf{Method} & \textbf{Domain}
  & \textbf{ECE}$\downarrow$ & \textbf{AdaECE}$\downarrow$ & \textbf{Brier}$\downarrow$
  & \textbf{AUROC}$\uparrow$ & \textbf{Acc} \\
\midrule
\multirow{4}{*}{Soft Verb.}
  & GSM8K      & .218 & .215 & .294 & .565 & .472 \\
  & MMLU       & .106 & .109 & .253 & .601 & .562 \\
  & ARC        & \textbf{.095} & .112 & .207 & .568 & .726 \\
  & TruthfulQA & .372 & .371 & .414 & .301 & .544 \\
\addlinespace
\multirow{4}{*}{P(True) Norm}
  & GSM8K      & .133 & .138 & .264 & .594 & .476 \\
  & MMLU       & .091 & .090 & .239 & .652 & .566 \\
  & ARC        & .117 & .132 & .210 & .624 & .728 \\
  & TruthfulQA & .089 & .092 & .180 & .818 & .532 \\
\addlinespace
\multirow{4}{*}{SECL}
  & GSM8K      & \textbf{.070} & \textbf{.079} & .249 & .573 & .488 \\
  & MMLU       & \textbf{.067} & \textbf{.060} & .249 & .549 & .558 \\
  & ARC        & .112 & \textbf{.106} & .207 & .616 & .714 \\
  & TruthfulQA & \textbf{.068} & .114 & .260 & .454 & .546 \\
\bottomrule
\end{tabular}
\end{table}

\begin{table}[h]
\centering
\caption{Per-domain results for Gemma~2-2B.
Best ECE per domain in \textbf{bold}.}
\label{tab:per_domain_gemma}
\small
\begin{tabular}{@{}llcccccc@{}}
\toprule
\textbf{Method} & \textbf{Domain}
  & \textbf{ECE}$\downarrow$ & \textbf{AdaECE}$\downarrow$ & \textbf{Brier}$\downarrow$
  & \textbf{AUROC}$\uparrow$ & \textbf{Acc} \\
\midrule
\multirow{4}{*}{Soft Verb.}
  & GSM8K      & .549 & .549 & .446 & .658 & .186 \\
  & MMLU       & .194 & .178 & .271 & .608 & .576 \\
  & ARC        & \textbf{.082} & \textbf{.085} & .201 & .534 & .738 \\
  & TruthfulQA & .267 & .263 & .338 & .417 & .566 \\
\addlinespace
\multirow{4}{*}{P(True) Norm}
  & GSM8K      & .395 & .394 & .326 & .685 & .186 \\
  & MMLU       & .163 & .175 & .263 & .616 & .576 \\
  & ARC        & .185 & .182 & .229 & .597 & .738 \\
  & TruthfulQA & \textbf{.104} & \textbf{.104} & .213 & .729 & .566 \\
\addlinespace
\multirow{4}{*}{SECL}
  & GSM8K      & \textbf{.356} & \textbf{.356} & .271 & .609 & .180 \\
  & MMLU       & \textbf{.054} & \textbf{.064} & .238 & .594 & .578 \\
  & ARC        & .144 & .136 & .213 & .591 & .728 \\
  & TruthfulQA & .210 & .198 & .293 & .343 & .574 \\
\bottomrule
\end{tabular}
\end{table}

\begin{table}[h]
\centering
\caption{Per-domain results for Phi~3.5-Mini.
Best ECE per domain in \textbf{bold}.}
\label{tab:per_domain_phi}
\small
\begin{tabular}{@{}llcccccc@{}}
\toprule
\textbf{Method} & \textbf{Domain}
  & \textbf{ECE}$\downarrow$ & \textbf{AdaECE}$\downarrow$ & \textbf{Brier}$\downarrow$
  & \textbf{AUROC}$\uparrow$ & \textbf{Acc} \\
\midrule
\multirow{4}{*}{Soft Verb.}
  & GSM8K      & .290 & .290 & .304 & .563 & .650 \\
  & MMLU       & .264 & .257 & .286 & .602 & .648 \\
  & ARC        & \textbf{.109} & \textbf{.105} & .150 & .588 & .826 \\
  & TruthfulQA & .343 & .336 & .362 & .592 & .546 \\
\addlinespace
\multirow{4}{*}{P(True) Norm}
  & GSM8K      & \textbf{.261} & \textbf{.257} & .303 & .559 & .650 \\
  & MMLU       & .171 & .158 & .243 & .648 & .648 \\
  & ARC        & .129 & .119 & .157 & .680 & .826 \\
  & TruthfulQA & \textbf{.146} & \textbf{.141} & .205 & .771 & .546 \\
\addlinespace
\multirow{4}{*}{SECL}
  & GSM8K      & .283 & .280 & .297 & .575 & .658 \\
  & MMLU       & \textbf{.054} & \textbf{.058} & .223 & .604 & .644 \\
  & ARC        & .129 & .129 & .167 & .499 & .826 \\
  & TruthfulQA & .229 & .224 & .319 & .407 & .532 \\
\bottomrule
\end{tabular}
\end{table}


% ============================================================================
\section{Ablation Studies} \label{app:ablations}
% ============================================================================

\paragraph{Weight accumulation.}
By default, SECL accumulates LoRA weights across the stream
(each question's adaptation builds on the previous state).
Table~\ref{tab:ablation_accum} shows that disabling accumulation
(resetting LoRA to the base model after each question)
is catastrophic: ECE degrades from 0.050 to 0.237, and AUROC
drops below chance (0.484), indicating that the model's confidence
signal is destroyed by isolated per-question updates.

\begin{table}[h]
\centering
\caption{Weight accumulation ablation (Llama~3.2-3B).}
\label{tab:ablation_accum}
\small
\begin{tabular}{@{}lcccc@{}}
\toprule
\textbf{Setting}
  & \textbf{ECE}$\downarrow$ & \textbf{Brier}$\downarrow$
  & \textbf{AUROC}$\uparrow$ & \textbf{Acc} \\
\midrule
SECL (accumulate)  & \textbf{.050} & \textbf{.241} & \textbf{.587} & .577 \\
Reset per question & .237 & .327 & .484 & .576 \\
\bottomrule
\end{tabular}
\end{table}

\paragraph{PH-gating vs.\ always-train.}
Training on every question (no PH-gating) achieves comparable ECE
to SECL (0.052 vs.\ 0.050), but at 4$\times$ the computational cost
since all 2{,}000 questions receive TTT updates.
PH-gating selects the 25.6\% of questions with the highest
entropy change, achieving the same calibration with 75\% fewer
gradient steps.

\begin{table}[h]
\centering
\caption{PH-gating vs.\ always-train (Llama~3.2-3B).
PH-gating matches always-train ECE at a fraction of the cost.}
\label{tab:ablation_phgate}
\small
\begin{tabular}{@{}lcccc@{}}
\toprule
\textbf{Setting}
  & \textbf{ECE}$\downarrow$ & \textbf{Brier}$\downarrow$
  & \textbf{AUROC}$\uparrow$ & \textbf{Trained \%} \\
\midrule
Always-train   & .052 & .242 & .585 & 100\% \\
PH-gated (SECL) & \textbf{.050} & \textbf{.241} & \textbf{.587} & 25.6\% \\
\bottomrule
\end{tabular}
\end{table}

\paragraph{KL divergence regularization.}
We experiment with adding a KL divergence term
$\beta \cdot D_{\text{KL}}(p_{\text{base}} \| p_{\text{adapted}})$
to the training loss to preserve the base model's output distribution.
Table~\ref{tab:ablation_kl} shows that a small $\beta{=}0.01$ slightly
improves Llama ECE (0.044 vs.\ 0.050) but degrades Gemma and Phi,
while $\beta{=}0.1$ is too aggressive for all models.
We use $\beta{=}0$ (no KL term) in all reported experiments.

\begin{table}[h]
\centering
\caption{KL regularization ablation across models (SECL + PH-gating).}
\label{tab:ablation_kl}
\small
\begin{tabular}{@{}llcccc@{}}
\toprule
\textbf{Model} & $\boldsymbol{\beta}$
  & \textbf{ECE}$\downarrow$ & \textbf{Brier}$\downarrow$
  & \textbf{AUROC}$\uparrow$ & \textbf{Acc} \\
\midrule
\multirow{3}{*}{Llama}
  & \textbf{0} (default) & .050 & .241 & \textbf{.587} & .577 \\
  & 0.01     & \textbf{.044} & \textbf{.242} & .591 & .573 \\
  & 0.1      & .149 & .279 & .527 & .569 \\
\addlinespace
\multirow{3}{*}{Gemma}
  & \textbf{0} (default) & \textbf{.056} & \textbf{.254} & \textbf{.548} & .515 \\
  & 0.01     & .127 & .271 & .486 & .514 \\
  & 0.1      & .238 & .305 & .551 & .520 \\
\addlinespace
\multirow{3}{*}{Phi}
  & \textbf{0} (default) & \textbf{.115} & \textbf{.251} & \textbf{.521} & .665 \\
  & 0.01     & .144 & .252 & .506 & .669 \\
  & 0.1      & .248 & .272 & .598 & .673 \\
\bottomrule
\end{tabular}
\end{table}


% ============================================================================
\section{Domain Order Sensitivity} \label{app:domain_order}
% ============================================================================

To assess sensitivity to the presentation order of domains, we run SECL 
with the reversed sequence 
(TruthfulQA $\to$ ARC $\to$ MMLU $\to$ GSM8K).
Table~\ref{tab:domain_order} compares overall metrics; 
Table~\ref{tab:domain_order_perdomain} gives per-domain detail.
SECL improves over Soft Verbalized under both orderings for all models,
though the original order yields lower ECE.
The gap is smallest for Llama (0.050 vs.\ 0.068), confirming that the 
method is reasonably robust to domain ordering on the strongest model.

\begin{table}[h]
\centering
\caption{Domain order sensitivity (overall metrics).
$\mathcal{O}$: GSM8K$\to$MMLU$\to$ARC$\to$TQA.
$\mathcal{R}$: TQA$\to$ARC$\to$MMLU$\to$GSM8K.}
\label{tab:domain_order}
\small
\begin{tabular}{@{}llccccc@{}}
\toprule
\textbf{Model} & \textbf{Order}
  & \textbf{ECE}$\downarrow$ & \textbf{AdaECE}$\downarrow$ & \textbf{Brier}$\downarrow$
  & \textbf{AUROC}$\uparrow$ & \textbf{Acc} \\
\midrule
\multirow{3}{*}{Llama 3.2-3B}
  & Soft Verb.       & .170 & .167 & .292 & .510 & .576 \\
  & SECL $\mathcal{O}$ & \textbf{.050} & \textbf{.060} & .241 & .587 & .577 \\
  & SECL $\mathcal{R}$ & .068 & .082 & .248 & .578 & .575 \\
\addlinespace
\multirow{3}{*}{Gemma 2-2B}
  & Soft Verb.       & .256 & .252 & .314 & .558 & .516 \\
  & SECL $\mathcal{O}$ & \textbf{.056} & \textbf{.060} & .254 & .548 & .515 \\
  & SECL $\mathcal{R}$ & .149 & .133 & .258 & .625 & .520 \\
\addlinespace
\multirow{3}{*}{Phi 3.5-Mini}
  & Soft Verb.       & .251 & .240 & .275 & .600 & .667 \\
  & SECL $\mathcal{O}$ & \textbf{.115} & \textbf{.119} & .251 & .521 & .665 \\
  & SECL $\mathcal{R}$ & .173 & .170 & .250 & .575 & .674 \\
\bottomrule
\end{tabular}
\end{table}

\begin{table}[h]
\centering
\caption{Per-domain ECE for original ($\mathcal{O}$) vs.\ reversed 
($\mathcal{R}$) domain order.}
\label{tab:domain_order_perdomain}
\small
\begin{tabular}{@{}llcccc@{}}
\toprule
\textbf{Model} & \textbf{Order}
  & \textbf{GSM8K} & \textbf{MMLU}
  & \textbf{ARC} & \textbf{TQA} \\
\midrule
\multirow{2}{*}{Llama}
  & $\mathcal{O}$ & .070 & .067 & .112 & .068 \\
  & $\mathcal{R}$ & .117 & .032 & .081 & .134 \\
\addlinespace
\multirow{2}{*}{Gemma}
  & $\mathcal{O}$ & .356 & .054 & .144 & .210 \\
  & $\mathcal{R}$ & .396 & .096 & .046 & .189 \\
\addlinespace
\multirow{2}{*}{Phi}
  & $\mathcal{O}$ & .283 & .054 & .129 & .229 \\
  & $\mathcal{R}$ & .068 & .213 & .099 & .343 \\
\bottomrule
\end{tabular}
\end{table}

\paragraph{ARC domain ordering effect.}
ARC-Challenge is the only domain where SECL consistently increases ECE
relative to the Soft Verbalized baseline in the forward order
(Llama: $.095 \to .112$; Gemma: $.082 \to .144$; Phi: $.109 \to .129$).
However, in the reversed order, ARC ECE \emph{improves} for all three models
(Llama: $.095 \to .081$; Gemma: $.082 \to .046$; Phi: $.109 \to .099$).
This confirms that the ARC degradation is an artifact of domain sequencing---specifically,
the LoRA weights adapted on preceding domains transfer poorly to ARC's
scientific reasoning format---rather than a fundamental limitation of SECL.
Overall ECE remains improved under both orderings for all models,
so the per-domain redistribution does not undermine the method's aggregate benefit.


% ============================================================================
\section{Computational Cost Analysis} \label{app:cost}
% ============================================================================

Table~\ref{tab:cost} compares the computational cost of SECL against
running P(True) Norm on every question.
We report costs in forward-pass equivalents (FWD-eq), where one 
backward pass $\approx$ 2 FWD.
For questions where PH-gating skips TTT, SECL requires only 1~FWD 
(the generation pass); the P(True) Norm baseline requires $1 + (k{+}1) = 6$~FWD 
per question ($k{=}4$ distractors plus the correct answer).

\begin{table}[h]
\centering
\caption{Computational cost comparison (in FWD-equivalents).
\emph{Trained} = questions receiving TTT;
\emph{Skipped} = PH-gated questions (generation only).
Per trained question: ${\approx}15$ FWD-eq 
(1~gen $+$ 5~P(True) $+$ 3~epochs $\times$ 3~FWD-eq/step).
Per skipped question: $1$ FWD-eq.
P(True) Norm baseline: 6 FWD-eq $\times$ 2{,}000 $= 12{,}000$.}
\label{tab:cost}
\small
\begin{tabular}{@{}lcccc@{}}
\toprule
\textbf{Model} & \textbf{Trained} & \textbf{Skipped}
  & \textbf{SECL cost} & \textbf{P(True) cost} \\
\midrule
Llama 3.2-3B & 512 (25.6\%) & 1{,}488 & 9{,}168 & 12{,}000 \\
Gemma 2-2B   & 119 ~~(5.9\%) & 1{,}881 & 3{,}666 & 12{,}000 \\
Phi 3.5-Mini & 160 ~~(8.0\%) & 1{,}840 & 4{,}240 & 12{,}000 \\
\bottomrule
\end{tabular}
\end{table}

SECL is cheaper than the P(True) Norm baseline for all three models.
Gemma and Phi trigger infrequently (6--8\%), reducing cost to roughly
a third of the baseline while achieving substantially better calibration.
Even Llama, which triggers most often (25.6\%), remains below the
baseline cost while improving ECE from .065 to .050.
Note that temperature scaling and Platt scaling add negligible compute
(a single parameter fit) but require ground-truth labels, which are 
unavailable in the test-time setting SECL targets.


% ============================================================================
\section{Negative Control: Qwen 2.5-3B} \label{app:qwen}
% ============================================================================

We evaluate Qwen~2.5-3B-Instruct as a negative control.
Unlike the other three models, Qwen's P(True) Norm baseline 
(\emph{best} $\tau{=}1.0$, ECE$=$0.257) is \emph{worse} than its
Soft Verbalized baseline (ECE$=$0.247), indicating that the 
discriminative signal provides no calibration advantage over generation.
Since SECL distills the discriminative signal into generative confidence,
the method cannot improve calibration when the discriminative signal 
itself is uninformative.
This confirms a key prerequisite of our approach: SECL succeeds 
precisely when a generation--discrimination gap exists.

\begin{table}[h]
\centering
\caption{Qwen~2.5-3B: no generation--discrimination gap.
All P(True) Norm temperatures yield worse ECE than Soft Verbalized.}
\label{tab:qwen}
\small
\begin{tabular}{@{}lc@{}}
\toprule
\textbf{Method} & \textbf{ECE}$\downarrow$ \\
\midrule
Soft Verbalized & \textbf{.247} \\
\addlinespace
P(True) Norm $\tau{=}0.3$ & .290 \\
P(True) Norm $\tau{=}0.7$ & .263 \\
P(True) Norm $\tau{=}1.0$ & .257 \\
P(True) Norm $\tau{=}1.5$ & .265 \\
P(True) Norm $\tau{=}2.0$ & .267 \\
P(True) Norm $\tau{=}3.0$ & .291 \\
\bottomrule
\end{tabular}
\end{table}


% ============================================================================
\section{PH-Gating Trigger Statistics} \label{app:ph_stats}
% ============================================================================

\begin{table}[h]
\centering
\caption{PH-gating trigger statistics per model.
\emph{Trained\%} is the fraction of questions receiving TTT adaptation.}
\label{tab:ph_triggers}
\small
\begin{tabular}{@{}lccc@{}}
\toprule
\textbf{Model} & \textbf{Total} & \textbf{Trained (\%)} & \textbf{Skipped (\%)} \\
\midrule
Llama 3.2-3B & 2{,}000 & 512 (25.6\%) & 1{,}488 (74.4\%) \\
Gemma 2-2B   & 2{,}000 & 119 ~~(5.9\%) & 1{,}881 (94.1\%) \\
Phi 3.5-Mini & 2{,}000 & 160 ~~(8.0\%) & 1{,}840 (92.0\%) \\
Llama 3.1-8B & 2{,}000 & 251 (12.6\%) & 1{,}749 (87.4\%) \\
\bottomrule
\end{tabular}
\end{table}


% ============================================================================
\section{Scaling to Llama 3.1-8B} \label{app:scaling_8b}
% ============================================================================

To test whether SECL scales to larger models, we evaluate
Llama~3.1-8B-Instruct (8.03B parameters) on the same 2{,}000-question stream.
We use LoRA on the last 8 layers (rank~8), $\tau{=}3.0$,
and the same PH-gating parameters as the 3B models.

\begin{table}[h]
\centering
\caption{Llama~3.1-8B results. SECL reduces ECE by 63\% relative
to Soft Verbalized while adapting only 12.6\% of questions.}
\label{tab:8b_results}
\small
\begin{tabular}{@{}lcccc@{}}
\toprule
\textbf{Method}
  & \textbf{ECE}$\downarrow$ & \textbf{Brier}$\downarrow$
  & \textbf{AUROC}$\uparrow$ & \textbf{Acc} \\
\midrule
Soft Verbalized      & .225 & .258 & \textbf{.684} & .644 \\
P(True) Norm         & .120 & \textbf{.211} & .718 & .646 \\
SECL (Ours)          & \textbf{.083} & .222 & .643 & .646 \\
\bottomrule
\end{tabular}
\end{table}

\begin{table}[h]
\centering
\caption{Per-domain results for Llama~3.1-8B.}
\label{tab:8b_per_domain}
\small
\begin{tabular}{@{}llcccc@{}}
\toprule
\textbf{Method} & \textbf{Domain}
  & \textbf{ECE}$\downarrow$ & \textbf{Brier}$\downarrow$
  & \textbf{AUROC}$\uparrow$ & \textbf{Acc} \\
\midrule
\multirow{4}{*}{Soft Verb.}
  & GSM8K      & .298 & .305 & .714 & .602 \\
  & MMLU       & .208 & .251 & .644 & .670 \\
  & ARC        & .108 & .164 & .633 & .798 \\
  & TruthfulQA & .287 & .310 & .669 & .506 \\
\addlinespace
\multirow{4}{*}{SECL}
  & GSM8K      & \textbf{.066} & .238 & .602 & .572 \\
  & MMLU       & \textbf{.120} & .223 & .591 & .686 \\
  & ARC        & .198 & .178 & .632 & .824 \\
  & TruthfulQA & \textbf{.038} & .249 & .552 & .504 \\
\bottomrule
\end{tabular}
\end{table}

SECL achieves ECE$=$0.083 on the 8B model, a 63\% reduction from the
Soft Verbalized baseline (0.225), confirming that the method scales
effectively. The pattern mirrors the 3B results: large improvements
on GSM8K, MMLU, and TruthfulQA, with ARC remaining resistant to
improvement (see domain ordering analysis in \S\ref{app:domain_order}).
The PH detector triggers on 12.6\% of questions, intermediate between
the 3B models' 5.9--25.6\% range.
